\chapter{Objetivos de Trabalho}
\section{Objetivo Geral}
Aprimorar e validar a arquitetura do IQArMobi, 
transformando - o em um agente inteligente para Monitoramento
da qualidade do ar, com capacidade de predição temporal e resposta automatizada
em tempo real, integrando a uma arquitetura IoT escalável baseada em sistemas
operacionais de tempo real e de middleware contexto, adaptando às particularidades
amazônicas.
\section{Objetivos Específicos}
\begin{enumerate}
    \item Expandir a arquitetura técnica, implementando um firware deterministico
          preemptivo baseado em FreeRTOS no microcontrolador Heltec Wifi LoRa 32 v2
          para gerenciar tarefas concorrente de leitura de sensores e comunicação via LoRaWAN
          \cite{alves2025freertos};
    \item Desenvolver e embarcar modelos de IA preditiva local 
          (regressão contínua de séries temporais) quantizados na 
          borda para antecipar concentrações quantitativas de gases
          poluentes (CO, O\textsubscript{3}, NH\textsubscript{3}, CO\textsubscript{2});
    \item Reescrever o midlware FIWARE para que ele se torne o Agente inteligente e ao mesmo tempo
          uma plataforma de gestão da qualidade do ar
          que vai garantir a interoperabilidade com outros sistemas de gestão da qualidade do ar
          em uma rede de sensores.  
    \item Validar o agente inteligente em campo por meio de uma campanha 
          experimental contínua de mais de 30 dias em 
          ambientes internos (laboratórios) e externos de um 
          smart campus na Amazônia, mensurando métricas de desempenho de 
          rede, consumo de hardware e precisão de IA
\end{enumerate}